% Further random-host refinement. Include after random-hosts-phase2.tex
% and k4-random-host-blocks.tex.
\section{A random-host upper coefficient of eight sevenths}
\label{sec:rhost-eight-sevenths}

\paragraph{Status.}
This section uses the explicit traceable $K_4$ blocks in Lemma~\ref{lem:rhost-k4-bands}.
It adds a uniform overlap estimate for their random embeddings.
This construction leaves the sharp coefficient and the full-threshold question unresolved.
Later sections improve the coefficient and address the threshold.

\begin{lemma}[Overlap estimate for a subdivided $K_4$]
\label{lem:rhost-k4-moments}
Let $Q$ be a subdivision of $K_4$, with $v$ vertices and $e=v+2$ edges.
Let $g$ be its shortest cycle length.
Let $b$ be the smallest number of edges in a connected subgraph of minimum degree at least two and rank two.
Let $J\sim G(n,D/n)$, with $D\ge54$ and $v^2/n=o(1)$.
Let $X$ count copies of $Q$, and let $Y$ count ordered intersecting pairs, including identical pairs.
Put $\mu=\mathbb E X$ and
\[
 \varepsilon_Q=\frac{v^2}{n}
 +(v+1)^{32}\left(D^{-g}+nD^{-b}+n^2D^{-e}\right).
\]
For an absolute constant $C$,
\[
 \operatorname{Var}(X)\le C\mu^2\varepsilon_Q,
 \qquad \mathbb E Y\le C\mu^2\varepsilon_Q.
\]
\end{lemma}

\begin{proof}
The normalized weight of an overlap with $u$ vertices and $f$ edges is
$(1+O(v^2/n))n^{-u}(n/D)^f$.
As in Lemma~\ref{lem:rhost-theta-moments}, all forest overlaps contribute $O(v^2/n)$.
Their disjoint components also supply at most the factor $\exp(2v^2/n)$ when a cyclic component is present.

A subdivision of $K_4$ has no two vertex-disjoint cycles.
Thus at most one overlap component contains a cycle.
Strip its degree-one vertices to obtain its nonempty core.
Each original subdivided edge is either wholly present in this core or absent from it.
The core is therefore one of the following:
a subdivided cycle, a subdivision of $K_4$ with one edge deleted, or all of $Q$.
Its excess $f-u$ is respectively zero, one, or two.
Its edge count is at least $g$, $b$, or $e$.

We give a polynomial bound for the number of connected cyclic overlap descriptions.
On each of the six original edge paths, such a component uses at most an initial segment from each endpoint.
Thus at most twelve cut positions specify its edge set, after the core type is chosen.
There are $O((v+1)^{12})$ possible source subgraphs.
Each has at most twelve vertices of degree one and four vertices of degree three.
Choose the images of these at most sixteen vertices in the second copy.
After suppressing degree-two vertices, at most twenty-four path segments remain.
Between two specified vertices in $Q$, there are at most $2^{\beta(Q)}=8$ simple paths.
A loop segment has at most fourteen oriented cycle images, since $Q$ has seven cycles.
Hence the segment images have only an absolute constant number of choices.
A component that is itself a cycle has at most $14v$ possible images and obeys the same bound.
There are therefore $O((v+1)^{28})$ connected cyclic overlap descriptions.

Their total normalized contribution is at most
\[
 O((v+1)^{28})\left(D^{-g}+nD^{-b}+n^2D^{-e}\right).
\]
The forest components and falling-factorial correction give the stated estimate with exponent $32$.
The argument applies first to labelled embeddings and then to unlabelled copies.
Vertex-disjoint pairs contribute at most $\mu^2$, which also proves the variance estimate.
\end{proof}

\begin{corollary}[A list of robust $K_4$ embeddings]
\label{cor:rhost-k4-list}
Let $d=(1+o(1))\ln n/\ln\ln n$ and $D=\Theta(d)$.
Assume $nD^{-d}$ tends to zero faster than every fixed negative power of $d$.
Consider $O(d)$ subdivisions of $K_4$.
Suppose each has $O(d)$ edges, shortest cycle length at least $d/4$,
every rank-two core has at least $d$ edges, and its full edge count is at least $2d$.
With high probability, all these graphs embed vertex-disjointly in $G(n,D/n)$.
They can avoid any prescribed set of $O(d^2)$ vertices, even if that set is chosen after observing the graph.
\end{corollary}

\begin{proof}
Apply Lemma~\ref{lem:rhost-k4-moments} to each graph in the list.
The sum of its error parameters, multiplied by $O(d^2)$, tends to zero.
The forest term contributes $O(d^5/n)$.
The cycle term tends to zero because $D^{-d/4}=n^{-1/4+o(1)}$.
The other terms tend to zero by the hypothesis on $nD^{-d}$.
The Chebyshev, Markov, and independent-set argument in Lemma~\ref{lem:rhost-theta-list}
therefore supplies more than $M$ vertex-disjoint copies of each type, for any fixed $M=O(d^2)$.
Choose one copy of each type greedily, avoiding the prescribed vertices and earlier choices.
\end{proof}

\begin{theorem}[An upper coefficient of eight sevenths]
\label{thm:rhost-eight-sevenths}
For the same $p_*$ as in Theorem~\ref{thm:rhost-scale}, and every $p\ge p_*$,
with high probability
\[
 \operatorname{Pex}(G(n,p))
 \le\log_2n+\left(\frac87+o(1)\right)\frac{\ln n}{\ln\ln n}.
\]
\end{theorem}

\begin{proof}
Use the exposure scheme and parameters of Theorem~\ref{thm:rhost-four-thirds}.
In particular, split $G_2$ into two independent subexposures of probability
$q_2=1-\sqrt{1-p_2}$, and use
\[
 \beta=\frac{4(\ln\ln n)^2}{\ln n},\qquad
 a=\lfloor(5\beta)^{-1}\rfloor,
 \qquad d=\lfloor\log_a(n/200)\rfloor.
\]
Let $b=\lceil\sqrt{2\ln n}\rceil$, $T=2(d+b+1)$, and $N=\lceil T/14\rceil$.
Replace the theta list by the two families of $K_4$ blocks from Lemma~\ref{lem:rhost-k4-bands}.
They use $2(N+1)$ blocks of rank three.
Add one single-cycle block for each length $3,\ldots,2N+5$, and one for length $4N+8$.
The resulting blocks cover every length through $14N+10\ge T$.
Their total rank is exactly
\[
 6(N+1)+(2N+4)=8N+10
 \le\frac87(d+b+1)+18.
\]

Prepare the single-cycle blocks, the binary base cycles, the two-digit base cycle,
and the reserve edge in $G_1$, using the separated subexposures described above.
Their total vertex count is $O(d^2)$.
Use the first $G_2$ subexposure for the complete $K_4$ blocks.
Corollary~\ref{cor:rhost-k4-list} applies.
Indeed, the metric lemma gives shortest cycles $2N+6>d/4$,
rank-two cores with at least $7N+13\ge d$ edges, and whole blocks with at least $14N+15\ge2d$ edges.
Also, $D=nq_2=\Theta(d)$ and $nD^{-d}$ has the required superpolynomial decay.

Each block has a specified unsubdivided arm and therefore has a spanning path.
Choose those paths after the embeddings are found.
Use the second $G_2$ subexposure to join the short blocks into the connector chain.
The previously audited endpoint-query and vertex-budget estimates remain unchanged.
Complete the binary and two-digit gadgets in $G_3$ and the Hamilton cycle in $G_4$.
Use the original $G_5$ long interval-covering chords.

The blocks cover every length through $T$.
The two-digit and binary gadgets cover the remaining intervals as before.
The final excess is at most
\[
 K+2b+8N+m+O(1)
 \le\log_2n+\frac87d+O(\sqrt{\ln n}).
\]
This gives the asserted asymptotic bound at $p_*$.
Monotonicity handles $p\ge p_*$.
\end{proof}

Together with Theorem~\ref{thm:rhost-coefficient-one}, this gives
\[
 \boxed{\quad
 \log_2n+(1-o(1))\frac{\ln n}{\ln\ln n}
 \le\operatorname{Pex}(G(n,p_*))
 \le\log_2n+\left(\frac87+o(1)\right)\frac{\ln n}{\ln\ln n}
 \quad}
\]
with high probability.
